%! Departures from Linearity — Unit 2, Lesson 2.5 — Student Packet Cover Sheet
\documentclass[12pt]{article}
\usepackage{apstats-article}
\usepackage{apstats-boxes}
\usepackage{ltablex}
\keepXColumns

\begin{document}

% Full-bleed navy banner. \coverbanner MEASURES the title block and sizes the
% band to it, so a lesson title long enough to wrap grows the band rather than
% running out the bottom of it.
\coverbanner{Unit 2}{Lesson 2.5 \quad Departures from Linearity}

% NAMESTRIP: this is the ONLY component that carries the name row.
\namedateperiod
\vspace{0.18in}

\boxguard[14]
\begin{learningtargetbox}
\small
% GRADUAL RELEASE: the vocabulary is taught in the guided notes, so the targets name it
% plainly. The AP Learning Objectives (DAT-1.I, DAT-1.J; DAT-1.F carried forward) stay in the
% teacher-facing plan.
\begin{enumerate}[leftmargin=1.5em, itemsep=5pt]
  \item \ldots{} tell an \textbf{outlier} from a \textbf{high-leverage point}, and decide from
        regression output whether a point is \textbf{influential}.
  \item \ldots{} use a \textbf{residual plot} to decide whether a line is an appropriate model
        --- even when $r$ is close to $1$ or $-1$.
  \item \ldots{} explain how a \textbf{transformation} such as taking the log of $y$ can
        straighten a curved pattern, and make a prediction from the transformed line.
\end{enumerate}
\end{learningtargetbox}

\vspace{0.14in}

\boxguard[14]
\begin{tocbox}
\small
\renewcommand{\arraystretch}{1.5}
\begin{tabularx}{\linewidth}{@{} c l X r @{}}
\rowcolor{sky}
\textbf{\#} & \textbf{Component} & \textbf{Description} & \textbf{Score} \\
\midrule
1 & Warm-Up & A taxi-fare residual, two residual plots, and one point that could move a line
                 & \blank{1.2cm} \\
2 & Guided Notes \& Practice & Three kinds of unusual point, a curve hiding behind
                 $r = -0.94$, and a lemonade stand's residual plot & \blank{1.2cm} \\
% AP Practice is assigned but NOT scored: its score cell prints NA, never a \blank{}.
% Row 4 is the lesson's GRADED practice. Paper homework is the DEFAULT for this course; on the
% occasions DeltaMath is assigned instead, that replaces row 4 and its score cell is struck.
3 & AP Practice & Used cars and a bacteria culture: AP-style multiple choice and one
                 free-response set --- practice, not a grade & \textbf{NA} \\
4 & Homework & Apartment rents and phone resale values: influence, residual plots, and a log
                 model --- \textbf{scored, due the first class after two study halls} & \blank{1.2cm} \\
\midrule
  & \multicolumn{2}{r}{\textbf{Total}} & \blank{1.2cm} \\
\end{tabularx}
\end{tocbox}

\vspace{0.14in}

\boxguard[8]
\begin{remindbox}
\small
A correlation close to $1$ or $-1$ says the points hug \emph{some} line tightly --- it does not
say a line is the right \emph{shape}. \textbf{Before you trust a line, look at its residual
plot}: random scatter means a line fits; a curve means it does not, whatever $r$ says. And
when one point sits apart from the rest, ask two separate questions: is it far from the
\emph{trend} (an outlier), and is it far out in $x$ (high leverage)? Then check what happens to
the slope and $r$ \textbf{without it} --- that is what makes a point influential.
\end{remindbox}

\end{document}
